\section*{LỜI CẢM ƠN}
\addcontentsline{toc}{section}{Lời cảm ơn}
Báo cáo được thực hiện trong khuôn khổ học phần Thuật toán ứng dụng, lớp N01, tại Đại học Phenikaa. Tập thể tác giả trân trọng cảm ơn giảng viên phụ trách học phần về các kiến thức nền tảng trong thiết kế thuật toán, phân tích độ phức tạp và mô hình hóa bài toán bằng đồ thị.

Quá trình nghiên cứu tạo điều kiện vận dụng các phương pháp đường đi ngắn nhất và mạng luồng vào bài toán lập lịch nhiều nguồn có ràng buộc. Việc kết hợp mô hình toán, cấu trúc dữ liệu và kiểm chứng thực nghiệm góp phần hình thành quy trình đánh giá thống nhất từ dữ liệu mặt bằng đến lịch sơ tán và các chỉ số hiệu quả.

Các góp ý học thuật về điều kiện biên, cách xác lập năng lực thông hành và phạm vi diễn giải kết quả có ý nghĩa quan trọng đối với việc hoàn thiện báo cáo. Những giới hạn liên quan đến dữ liệu hiện trường, diễn biến nguy cơ và hành vi di chuyển được trình bày nhằm xác định rõ phạm vi áp dụng của mô hình.

Mặc dù đã thực hiện kiểm thử thuật toán và đối chiếu lịch sau tối ưu, nghiên cứu vẫn còn những vấn đề cần khảo sát thêm về hiệu chỉnh tham số và đánh giá độ nhạy. Các nhận xét chuyên môn là cơ sở để tiếp tục hoàn thiện phương pháp và mở rộng thực nghiệm.

\begin{flushright}\textit{Tập thể tác giả}\\Lý Thành Đạt\\Đinh Đức Mạnh\\Phạm Văn Đạt\end{flushright}
